% !TeX root = ../main.tex
\section{Simply typed lambda calculus}

The untyped lambda calculus of \S2 gives us a vehicle for computation, but it's a vehicle with no guardrails: $x\,x$, $(\lambda x.\,x)\,(\lambda y.\,y\,y)$, and the Y combinator are all perfectly legal terms, and nothing in the syntax stops you from writing $\mathit{TRUE}\;3$ or applying a Church numeral to a boolean. The \textbf{simply typed lambda calculus} (STLC) is the untyped calculus plus the smallest possible type discipline layered on top: every term gets classified by a type, and application is only legal when the types line up. It's the smallest interesting typed language, and nearly everything in the rest of this chapter is STLC plus one more feature.

\subsection{types and terms}

We add a syntactic category of \textbf{types}, built from a couple of base types and a function-type former:
\[
  \tau ::= \tyint \mid \tybool \mid \tau \to \tau
\]
$\tyint$ and $\tybool$ classify integer and boolean literals; $\tau_1 \to \tau_2$ classifies functions that take a $\tau_1$ and return a $\tau_2$. (Function types associate to the right, so $\tyint \to \tyint \to \tybool$ means $\tyint \to (\tyint \to \tybool)$ — the type of a two-argument curried function.)

Terms are the same shape as before, except abstractions now carry an explicit annotation on their parameter:
\[
  e ::= x \mid \lambda x{:}\tau.\, e \mid e\, e \mid n \mid \ift{e}{e}{e}
\]
where $n$ ranges over integer literals. We've thrown in literals and $\ift{\cdot}{\cdot}{\cdot}$ because a type system with nothing to disagree about (just pure functions) is a little unconvincing — you want at least two base types so that ``ill-typed'' is a real possibility, not a vacuous one.

\subsection{the typing judgment}

The central object of study is the judgment $\Gamma \vdash e : \tau$: under context $\Gamma$ (a finite map from variables to types), term $e$ has type $\tau$. The rules are syntax-directed — one rule per term constructor — which means that, reading bottom-up, there is exactly one rule that could possibly apply to any given expression. That's what makes type \emph{checking} for STLC mechanical: no guessing.

\[
\infer[\textsc{T-Var}]
  {\Gamma \vdash x : \tau}
  {x : \tau \in \Gamma}
\qquad
\infer[\textsc{T-Abs}]
  {\Gamma \vdash \lambda x{:}\tau_1.\, e : \tau_1 \to \tau_2}
  {\Gamma, x : \tau_1 \vdash e : \tau_2}
\qquad
\infer[\textsc{T-App}]
  {\Gamma \vdash e_1\, e_2 : \tau_2}
  {\Gamma \vdash e_1 : \tau_1 \to \tau_2 \quad \Gamma \vdash e_2 : \tau_1}
\]

\textsc{T-Var} says a variable's type is whatever $\Gamma$ says it is. \textsc{T-Abs} says a function's type is $\tau_1 \to \tau_2$ exactly when its body has type $\tau_2$ under the extra assumption that the parameter has type $\tau_1$ — note the annotation on the lambda is what tells us $\tau_1$; nothing here is being guessed. \textsc{T-App} is the rule that actually enforces the discipline: to apply $e_1$ to $e_2$, $e_1$'s type must literally be a function type whose domain matches $e_2$'s type. This is exactly the rule that rejects $3\,(\lambda x.\,x)$ ($3$ isn't a function type at all) and $(\lambda x{:}\tyint.\,x)\,\mathit{true}$ (domain mismatch: $\tyint \neq \tybool$).

We also add the obvious rules for literals and conditionals, which we won't dwell on:
\[
\infer[\textsc{T-If}]
  {\Gamma \vdash \ift{e_1}{e_2}{e_3} : \tau}
  {\Gamma \vdash e_1 : \tybool \quad \Gamma \vdash e_2 : \tau \quad \Gamma \vdash e_3 : \tau}
\]

\subsection{capture-avoiding substitution}

$\beta$-reduction, recall from \S2, is $(\lambda x.\,e)\,v \to e[x := v]$. We've been writing $e[x := v]$ as if ``substitute $v$ for $x$ in $e$'' were self-explanatory, but there's a subtlety hiding in it: substitution must be \textbf{capture-avoiding}. Naively replacing every free occurrence of $x$ by $v$ can accidentally trap a free variable of $v$ inside a binder in $e$ that happens to reuse the same name.

\begin{example}[why naive substitution breaks]
Take $e = \lambda y.\, x\, y$ and substitute $v = y$ (a variable named $y$, free) for $x$. Naive, textual substitution gives $\lambda y.\, y\, y$ — but that's wrong: the $y$ we substituted in was free before the substitution, and now it's been captured by the $\lambda y$ that happened to already be there. The two occurrences of $y$ in the result mean completely different things, even though they're spelled the same.

The fix uses $\alpha$-equivalence (\S2): before substituting, rename the bound variable $y$ in $e$ to something fresh, say $z$, giving $\lambda z.\, x\, z \equiv_\alpha \lambda y.\, x\, y$. \emph{Now} substitute: $\lambda z.\, y\, z$. The free $y$ stays free, exactly as it should.
\end{example}

In short, capture-avoiding substitution is ``substitute, but rename bound variables in $e$ first if they'd otherwise collide with a free variable of $v$''. It's a nuisance to implement (this is one of the reasons why real compilers often use de Bruijn indices instead of named variables internally), but conceptually it's just $\alpha$-equivalence doing its job: since bound-variable names are meaningless, we're always free to pick fresh ones.

\subsection{strong normalization}

Here is the first serious payoff of adding types. Recall the Y combinator from \S2, $Y = \lambda f.\,(\lambda x.\, f\,(x\, x))\,(\lambda x.\, f\,(x\, x))$, which lets untyped lambda calculus express unbounded recursion — and, as a consequence, lets it diverge: $Y\,g$ never reaches a normal form for most $g$, by design. STLC cannot do this.

\begin{theorem}[strong normalization]
If $\vdash e : \tau$ for some type $\tau$ (i.e., $e$ is well-typed in the empty context), then every reduction sequence starting from $e$ terminates in a normal form. There is no well-typed term in STLC that diverges.
\end{theorem}

This is a genuinely striking fact: it means STLC, unlike the language it's built on top of, is not Turing-complete. Every well-typed STLC program halts. (This is also why real programming languages need \emph{more} than STLC — recursion is useful! We'll see in \S7 how recursive types let you smuggle recursion back in deliberately, rather than get it by accident.)

\emph{Why can't the Y combinator be typed?} Try it. $Y$'s definition hinges on the self-application $x\, x$ inside $\lambda x.\, f\,(x\,x)$. For $x\,x$ to typecheck, $x$ needs a function type whose domain is $x$'s own type — say $x : \alpha$, and then \textsc{T-App} demands $\alpha = \alpha \to \beta$ for some $\beta$. But $\tau ::= \tyint \mid \tybool \mid \tau \to \tau$ only contains \emph{finite} trees: there is no finite type $\alpha$ satisfying $\alpha = \alpha \to \beta$, because the right-hand side is always one constructor taller than the left. (We'll meet this exact equation again in \S4, as the reason Hindley-Milner's occurs check rejects $\lambda f.\,f\,f$, and again in \S7, where recursive types $\mu\alpha.\,\tau$ are introduced specifically to give equations like this a solution.) So self-application is simply not well-typed in STLC — no typing derivation for $x\,x$ exists in any context — and since $Y$ is built entirely out of self-applications, $Y$ has no type either. No self-application means no unbounded unfolding, which is the deep reason strong normalization holds.

\subsection{Curry-Howard: propositions as types}

Here's a beautiful accident (or non-accident) of STLC's design: its types, read as logical formulas, are exactly the propositions of intuitionistic propositional logic, and its well-typed terms are exactly the proofs of those propositions. This is the \textbf{Curry-Howard correspondence}, and the dictionary looks like this:

\begin{center}
\begin{tabular}{ll}
  \textbf{Logic} & \textbf{Types} \\
  \hline
  Proposition & Type \\
  Proof of $P$ & Term of type $P$ \\
  Implication $P \Rightarrow Q$ & Function type $P \to Q$ \\
  Conjunction $P \wedge Q$ & Product/pair type $P \times Q$ \\
  Modus ponens & Function application \\
\end{tabular}
\end{center}

Under this reading, \textsc{T-App} \emph{is} modus ponens: from a proof of $P \Rightarrow Q$ (a function $e_1 : \tau_1 \to \tau_2$) and a proof of $P$ (a term $e_2 : \tau_1$), you get a proof of $Q$ (the term $e_1\,e_2 : \tau_2$). And \textsc{T-Abs} is exactly conditional proof: to prove $P \Rightarrow Q$, assume $P$ and derive $Q$ — which is precisely ``extend $\Gamma$ with $x : \tau_1$ and typecheck the body against $\tau_2$''.

\begin{example}[a small proof-as-program]
The proposition ``$P \wedge Q \Rightarrow P$'' (and-elimination) corresponds to the type $\tau_1 \times \tau_2 \to \tau_1$. Its proof is the program that projects the first component of a pair:
\[
  \lambda p{:}\tau_1\times\tau_2.\; \mathsf{fst}(p)
\]
Typechecking this term \emph{is} verifying the proof. Conversely, the fact that $\tau_1 \to \tau_2 \to \tau_1$ (i.e. $P \Rightarrow (Q \Rightarrow P)$, a logical tautology: ``if $P$, then anything implies $P$'') is inhabited by the term $\lambda x{:}\tau_1.\,\lambda y{:}\tau_2.\,x$ is itself a proof that the tautology holds. Finding a well-typed term of a given type and finding a proof of a given proposition are, quite literally, the same search problem.
\end{example}

This correspondence is more than a curiosity — it's the theoretical basis for proof assistants like Coq and Agda, where ``writing a program'' and ``proving a theorem'' are the same activity, checked by the same type checker.

\subsection{Exercises}
\begin{itemize}
  \item Give a full typing derivation (using \textsc{T-Var}, \textsc{T-Abs}, \textsc{T-App}) for $\vdash \lambda x{:}\tyint.\,\lambda y{:}\tyint.\, x : \tyint \to \tyint \to \tyint$.
  \item Is $\lambda x{:}(\tyint\to\tyint).\, x\, x$ well-typed in STLC? Explain precisely which rule fails and why, connecting your answer to the strong normalization argument above.
  \item Perform the substitution $(\lambda y.\, x\,(\lambda x.\, x\,y))[x := y\,z]$ correctly (i.e., capture-avoiding), showing which bound variable needs renaming and why.
  \item Under Curry-Howard, what proposition does the type $(\tau_1 \to \tau_2) \to (\tau_2 \to \tau_3) \to (\tau_1 \to \tau_3)$ correspond to? Write down a well-typed term of that type (it's function composition) and identify which logical rule it witnesses.
  \item STLC is not Turing-complete, yet real typed languages obviously are. Name one language feature (not covered yet in this chapter) whose whole purpose is to reintroduce unbounded recursion into an otherwise strongly-normalizing type system.
\end{itemize}
